\section{When the Lubin--Tate branch bound is sharp}
\label{sec:branch-equality}

This calculation answers the question about cancellation left by the
\href{https://github.com/VladimirReshetnikov/ProveIt/blob/5cc267af75a635199ac5d93a7d49cbce521daefc/docs/incoming/grh_transfers_cycle_signings_and_lubin_tate.zip}{earlier branch-interpolation report}.
Its source in the Lean study was the use of an extremizer to establish the
sharpness of a derivative bound
(\href{https://github.com/lsernau/grh-lean/blob/72fead102f663d5c7a407964a13dc0cf9d977207/grh/src/New/Derivatives/Sharpness/OneCurve.lean\#L165-L183}{the one-curve lemma}).
The application concerns finite precision for the branch-evaluation method
in Berger and Sprang's \href{https://arxiv.org/abs/2502.18053}{Lemma~4.2 and Proposition~4.4}.

Let $\mathcal O$ be the ring of integers of a finite extension of
$\mathbb Q_p$, with uniformizer $\pi$, valuation $v(\pi)=1$ and residue
field of size $q$. Fix a Lubin--Tate formal group for $\pi$.
Its endomorphism $[a](X)$ has linear term $aX$ and
$[\pi](X)\equiv X^q\pmod\pi$ in a Lubin--Tate coordinate.
Let $A=\mathcal O/\pi^r$, $r\ge1$.
For a finite set $B\subset\mathcal O$ of at least two distinct directions
and $b\in B$, put
\[
 h_c=v(b-c),\quad h=\max_{c\ne b}h_c,\qquad
 \omega=\sum_{c\ne b}q^{h_c},\qquad
 \Delta=\prod_{c\ne b}([b](X)-[c](X)).
\]
Write $t$ for the least integer such that $X^t/\Delta\in A[[X]]$.
The previous report bounds $t\le\omega+(r-1)(q^h-1)$.
The branch cutoff is $1+t$ when recovering all coefficients of total degree
less than $|B|$. Indeed the Lagrange numerator
$P_b=\prod_{c\ne b}(Y-[c](X))$ has total order $|B|-1$ and leading
$Y$-coefficient one. If all observations vanish at their indicated
precisions, each term $H(X,[b](X))P_b/\Delta$ of the interpolation
remainder has total order at least $|B|$. The discarded multiple of
$\prod_c(Y-[c](X))$ has the same property. Division by this monic
distinguished polynomial is valid over the complete local ring $A[[X]]$.
Conversely $(X^t/\Delta)P_b$ has a nonzero constant-in-$X$ coefficient of $Y^{|B|-1}$
and vanishes on every other branch. Its value on branch $b$ is $X^t$.

\paragraph{Equality criterion.}
The bound is an equality if $r=1$ or $h=0$. For $r>1$ and $h>1$ it is
always strict. For $h=1$, set $N=\#\{c\ne b:h_c=1\}$ and define
\[
 j_* =\max\left\{0\le j<r:
         j+v\!\binom{N+j-1}{j}<r\right\}.
\]
Here $v$ acts on the indicated positive integer, with the normalization
above. The exact pole order is
\[
                    t=\omega+(q-1)j_* .
\]
Consequently at height one equality in the earlier bound holds exactly
when $p$ does not divide $\binom{N+r-2}{r-1}$.

\paragraph{Proof.}
The difference $[b]-[c]$ is $[\pi^{h_c}]$ times a unit in $\mathcal O[[X]]$.
In fact $F(U,V)-U$ equals $V$ times a unit, where $F$ is the formal
addition law. Substitute $U=[c]$ and $V=[b-c]$, and use
$[b-c]=[u]\circ[\pi^{h_c}]$ with $u$ a unit.
Multiplication by a power-series unit preserves the least exponent with
nonzero coefficient in $A((X))$, including when that coefficient is a
zero divisor. Hence all differences at the same height can be replaced
by identical factors.

Height zero gives $[1](X)=X$. For $a\ge1$, put $s=q^a$.
Since $[\pi^a](X)\equiv X^s\pmod\pi$ and the linear term is
$\pi^aX$, write
\[
            [\pi^a](X)=X^s U_a(X)+\pi X V_a(X),\qquad
            U_a(0)\in\mathcal O^\times,\quad V_a(0)=\pi^{a-1}.
\]
Take the terms of degrees below $s$ into $\pi XV_a$.
Inverting the factor in $A((X))$ amounts to expanding
$(1+\pi X^{1-s} V_a/U_a)^{-1}$ through power $r-1$ and multiplying
by $X^{-s}U_a^{-1}$. In the product of these expansions at most
$r-1$ powers of $\pi$ can contribute. Each can add at most $q^h-1$
to the pole order. This gives the earlier bound. Modulo $\pi$ the
reciprocal has pole order $\omega$, which proves equality if $r=1$.
If $h=0$, all factors are $X$ times units at every precision.

Suppose $h>1$ and $r>1$. A contribution at the extreme exponent would
require $r-1$ factors of height $h$, counted with repetition, and would
use only their constant terms $V_h(0)/U_h(0)$. Each such constant is
divisible by $\pi$. The coefficient therefore vanishes modulo $\pi^r$.

Finally suppose $h=1$. Up to a power-series unit the reciprocal is
\[
 X^{-\omega}
 \sum_{j=0}^{r-1}(-1)^j\binom{N+j-1}{j}
       \pi^j X^{-j(q-1)}(V_1/U_1)^j .
\]
The power series $V_1/U_1$ is a unit. Each summand whose scalar survives has
least exponent $-\omega-j(q-1)$. The largest surviving $j$ thus
determines the least exponent in the whole sum. This proves the formula.

\paragraph{Exact answer modulo $\pi^2$ at arbitrary contact height.}
For each $a\ge1$ put $N_a=\#\{c\ne b:h_c=a\}$.
If all $N_a$ are divisible by $p$, then $t=\omega$ at $r=2$.
Otherwise let $a_*$ be the largest $a\ge1$ for which $p\nmid N_a$. Then
\[
                  t=\omega+(q-1)q^{a_*-1}\qquad(r=2).
\]
To see this, write $[\pi](X)=X^q+\pi W(X)$, where $W(0)=0$ and
$W'(0)=1$. Reduction of $[\pi^a]=[\pi]\circ[\pi^{a-1}]$ gives
$[\pi^a](X)\equiv X^{q^a}+\pi W(X^{q^{a-1}})\pmod{\pi^2}$:
adding a multiple of $\pi$ to the argument leaves its $q$th power
unchanged modulo $\pi^2$, since $p\mid q$. Up to a power-series unit,
the reciprocal of $\Delta$ is therefore
\[
 X^{-\omega}\left(1-
 \pi\sum_{a\ge1}N_a X^{-q^a}W(X^{q^{a-1}})\right)\pmod{\pi^2}.
\]
Each surviving term indexed by $a$ has its first power at
$-(q-1)q^{a-1}$ inside the parentheses. These exponents are distinct,
and the largest surviving height determines $t$.

For the multiplicative formal group over $\mathbb Z_2$ take $r=2$.
At $b=0$ for $B=\{0,2,6\}$ one has $N=2$, $\omega=4$ and $j_*=0$,
so the required cutoff is $5$; the former bound gave $6$.
For $B=\{0,4\}$ at $b=0$, the polynomial
$[4]=(1+X)^4-1$ reduces to $X^4+2X^2$ modulo $4$.
Its reciprocal is $X^{-4}-2X^{-6}$, giving $t=6$ against the bound $7$.
